\documentclass[10pt]{article}

% --- Page Setup ---
% \usepackage[a4paper, margin=0.5cm]{geometry}
\usepackage[a4paper, left=0.5cm, right=0.5cm, top=0.3cm, bottom=0.3cm]{geometry}
\usepackage{multicol}

% \allowdisplaybreaks

% --- Math and Code ---
\usepackage{amsmath, amssymb}
\usepackage{listings}

% --- Spacing Adjustments ---
\usepackage{enumitem}
\usepackage{titlesec}

\usepackage{nicematrix, tikz}

% 1. Remove space between list items
\setlist{nosep, leftmargin=*} 

% 2. Shrink space around section headers
\titlespacing*{\section}{0pt}{0pt}{0pt}
\titlespacing*{\subsection}{0pt}{1ex plus .2ex minus .2ex}{0.5ex plus .1ex}

% --- NEW: Reduce header font sizes ---
\titleformat*{\section}{\normalsize\bfseries\color{blue!75!black}}
\titleformat*{\subsection}{\small\bfseries\itshape}

% 3. Add a thin line between columns for readability
\setlength{\columnseprule}{0.4pt}
\setlength{\columnsep}{1.5em}

% --- Custom Code Formatting ---
\lstset{
    basicstyle=\ttfamily\scriptsize,
    breaklines=true,
    aboveskip=2pt,
    belowskip=2pt
}

\usepackage{xcolor}
% \newcommand{\sectionbox}[1]{%
%      \colorbox{blue!75!black}{\color{white}\small\bfseries~#1~}%
%    }
%    \titleformat{\section}{}{}{0em}{\sectionbox}
%    \titleformat*{\subsection}{\normalsize\bfseries}

\setlength{\parskip}{0pt}
\setlength{\baselineskip}{0pt plus 0.2pt} 

\begin{document}
\small 

\raggedright
% \raggedbottom

\begin{multicols*}{2}

\section*{General stuff}
$\|\mathbf{u} + \mathbf{v}\|^2 + \|\mathbf{u} - \mathbf{v}\|^2 = 2(\|\mathbf{u}\|^2 + \|\mathbf{v}\|^2)$.\\
$\begin{bmatrix}a & b \\ c & d\end{bmatrix}^{-1} = \frac{1}{ad-bc}\begin{bmatrix}d & -b \\ -c & a\end{bmatrix}, \det\begin{bmatrix} 
a & b & c \\ 
d & e & f \\ 
g & h & i 
\end{bmatrix} = a(ei - fh) - b(di - fg) + c(dh - eg)$\\
$\det(A^T) = \det(A)$.\\
$\sin(\theta) = \frac{e^{i\theta} - e^{-i\theta}}{2i}$, $\cos(\theta) = \frac{e^{i\theta} + e^{-i\theta}}{2}$\\
$\sin(\alpha \pm \beta) = \sin(\alpha)\cos(\beta) \pm \cos(\alpha)\sin(\beta)$, $\cos(\alpha \pm \beta) = \cos(\alpha)\cos(\beta) \mp \sin(\alpha)\sin(\beta)$

\section*{Vectors}
$\mathbf{u}\cdot\mathbf{v}=|\mathbf{u}||\mathbf{v}|\cos\theta = u_1v_1 + \ldots + u_nv_n$. $\theta$ acute $\iff \mathbf{u}\cdot \mathbf{v}>0$.\\
\textbf{Cauchy-Schwarz Ineq}: $|\mathbf{u} \cdot \mathbf{v}| \le ||\mathbf{u}|| ||\mathbf{v}||\to -1\leq\frac{\mathbf{u}\cdot\mathbf{v}}{||\mathbf{u}||||\mathbf{v}||}\leq 1$\\
\textbf{Triangle Ineq}: $||\mathbf{u}+\mathbf{v}|| \le ||\mathbf{u}|| + ||\mathbf{v}||$\\
\textbf{Vector/Parametric Eqn of $k$-D object in $\mathbb{R}^n$}: $\mathbf{x} = \mathbf{x}_0 + t_1\mathbf{v}_1 + t_2\mathbf{v}_2 + \dots + t_k\mathbf{v}_k$\\
\textbf{Point-normal Eqn}: $\mathbf{n} \cdot (\mathbf{x} - \mathbf{x}_0) = 0$ describes a hyperplane, an object with exactly one dimension less than the surrounding space.

\section*{Matrices}
trace($A$) $= a_{11} + a_{22} + \ldots + a_{nn}$.\\
Upper triangular: all entries below main diagonal are zeroes.\\
Diagonal: all entries outside the main diagonal are zeroes.\\
\textbf{Symmetric}: $A^T = A$.\\
\textbf{Skew-symmetric}: $A^T = -A$. Diagonal all 0s.\\
\textbf{Transpose}: $(A^T)^T=A$, $(A+B)^T = A^T + B^T$, $(AB)^T = B^TA^T$.\\
\textbf{Inverse}: $AA^{-1} = A^{-1}A = I$. $(AB)^{-1} = B^{-1}A^{-1}$, $(A^T)^{-1} = (A^{-1})^T$.\\
\textbf{Multiplicatn as linear comb}: $A\mathbf{x} = x_1\begin{bmatrix} a_{11} \\ a_{21} \\ \vdots \\ a_{m1} \end{bmatrix} + x_2\begin{bmatrix} a_{12} \\ a_{22} \\ \vdots \\ a_{m2} \end{bmatrix} + \dots + x_n\begin{bmatrix} a_{1n} \\ a_{2n} \\ \vdots \\ a_{mn} \end{bmatrix}$\\
\textbf{Row method}: $AB = \begin{bmatrix} \mathbf{a}_1 \\ \mathbf{a}_2 \\ \vdots \\ \mathbf{a}_m \end{bmatrix} B = \begin{bmatrix} \mathbf{a}_1 B \\ \mathbf{a}_2 B \\ \vdots \\ \mathbf{a}_m B \end{bmatrix}$\\
\textbf{Col method}: $AB = A[\mathbf{b}_1 \quad \mathbf{b}_2 \quad \dots \quad \mathbf{b}_n] = [A\mathbf{b}_1 \quad A\mathbf{b}_2 \quad \dots \quad A\mathbf{b}_n]$\\

\section*{SLE}
For linear equation in $n$ vectors: $x_1\textbf{v}_1 + \ldots x_n\mathbf{v}_n = \mathbf{b}$,
\begin{itemize}
    \item Span$\{\mathbf{v}_1\ldots\mathbf{v}_n\}$ is set of all linear combinations of $\mathbf{v}_1\ldots\mathbf{v}_n$.
    \item The homogenous eqn ($\mathbf{b}=\mathbf{0}$) is always consistent.
    \begin{itemize}
        \item $\mathbf{v}_1 \ldots \mathbf{v}_n$ are LI $\Leftrightarrow$ there is only the trivial solution ($x_i=0$).
        \item Otherwise, $\exists$ non-zero $x_1\ldots x_n$ that satisfies the homo eqn.
        \item If one of the vectors is $\mathbf{0}$, the set is LD.
    \end{itemize}
\end{itemize}
\textbf{Row reduction} Row equivalent matrices can be transformed to eo by EROS:
\begin{enumerate}
    \item Multiply row by a \textbf{non-zero} constant.
    \item Swap two rows.
    \item Add multiple of one row to another.
\end{enumerate}
EROs manipulate the hyperplane without affecting the intersections, until they are parallel to the standard $x, y, z$ coordinate axes.\\
\textbf{REF (Not unique)}
\begin{itemize}
    \item All zero rows at the bottom.
    \item All pivots (first non-zero entry in a row) is strictly to the right of pivot above it.
    \item \textbf{RREF (Unique)}: Pivots = 1, and every other number in that column is 0.
\end{itemize}
\textbf{Gaussian elim}: forward phase to get REF.\\
\textbf{Gauss-Jordan elim}: forward + backward phase to get RREF.\\
% \textbf{Gauss and Gauss-Jordan elimination} use EROs to achieve REF and RREF.\\
Interpreting REF/RREF form of augmented matrix $[A \mid \mathbf{b}]$:
\begin{enumerate}
    \item If a row $[0 \quad 0 \quad 0 \mid c], c \neq 0$ appears, system is \textbf{inconsistent}. $\mathbf{b}$ is out of Span ($A$).
    \item Else, continue to RREF or write down the eqns.
    \item Express each leading variable in terms of free variables.
    \item \textbf{Unique soln}: Every col of $A$ has a pivot. Cols of $A$ LI.
    \item \textbf{Infinite solns}: $\exists$ non-pivot cols, hence $\exists$ free variables. Cols of $A$ LD.
\end{enumerate}
\textbf{Pivots}
\begin{itemize}
    \item Pivot in every col $\Leftrightarrow$ no free variables $\Leftrightarrow$ $\leq 1$ soln $\Leftrightarrow$ LI cols. $A\mathbf{x}=\mathbf{0}$ has only trivial soln.
    \item Pivot in every row $\Leftrightarrow$ cols span entirety of $\mathbb{R}^m$ $\Leftrightarrow$ $\exists$ solution for every $\mathbf{b} \in \mathbb{R}^m$.
    \item Pivot in every row and col $\Leftrightarrow$ square matrix $\Leftrightarrow$ unique soln for every $\mathbf{b}$ $\Leftrightarrow$ $A$ invertible and $\mathbf{x} = A^{-1}\mathbf{b}$.
    \item Wide matrices: cannot have pivot in every col $\Leftrightarrow$ alw have free variables $\Leftrightarrow$ either inconsistent or infinite solns.
    \item Tall matrices: cannot have pivot in every row $\Leftrightarrow$ alw have zero row $\Leftrightarrow$ no soln for some $\mathbf{b}$.
\end{itemize}
\textbf{Linear Transformation} $T(\mathbf{x}) = A\mathbf{x}$, where
\begin{itemize}
    \item $T(\mathbf{u} + \mathbf{v}) = T(\mathbf{u}) + T(\mathbf{v})$
    \item $T(c\mathbf{u}) = cT(\mathbf{u})$
\end{itemize}
Unique $A$ can constructed by transforming each basis vector, i.e. $A = [T(\mathbf{e}_1) \quad T(\mathbf{e}_2) \quad \dots \quad T(\mathbf{e}_n)]$\\
For linear transform of matrices, flatten each matrix to a basis vector.

\section*{Invertibility}
\textbf{Elementary matrix}: Created by performing 1 ERO on $I$.\\
Left-multiplying $A$ by $E$ ($EA$) applies that ERO on $A$.\\
Every elem matrix $E$ is invertible. $E^{-1}$ reverses the same operation.\\
Deriving $A^{-1}$ by Gauss-Jordan: $[A \mid I]$ $\xrightarrow{\text{row reduce}}$ $[I \mid A^{-1}]$.\\
\textbf{Invertible Matrix Theorem} For $n \times n$ $A$,
\begin{enumerate}
    \item $A$ invertible.
    \item $A$ row equivalent to $I_n$.
    \item $A$ has $n$ pivots.
    \item $A\mathbf{x}=\mathbf{0}$ has only the trivial soln.
    \item Cols of $A$ LI.
    \item $A\mathbf{x}=\mathbf{b}$ has unique soln for every $\mathbf{b}\in\mathbb{R}^n$.
    \item Cols of $A$ span $\mathbb{R}^n$.
    \item $\exists$ matrix $C$ s.t. $CA=I$ (left inverse).
    \item $\exists$ matrix $D$ s.t. $AD=I$ (right inverse).
    \item $A^T$ is invertible.
    \item Cols of $A$ form basis for $\mathbb{R}^n$.
    \item $\mathcal{C}(A)=\mathbb{R}^n$.
    \item $\dim \mathcal{C}(A) = n$.
    \item rank($A$) = $n$.
    \item $\mathcal{N}(A) = \{\mathbf{0}\}$.
    \item $\dim \mathcal{N}(A) = 0$.
\end{enumerate}
\textbf{LU factorisation} For solving sequence of eqns like $A\mathbf{x}=\mathbf{b}_1$, $\ldots$ $A\mathbf{x}=\mathbf{b}_p$, LU fact is faster than computing $A^{-1}$.\\
Decompose $A_{m\times n} = L_{m \times m}U_{m \times n}$, where $L$ is \textbf{unit} lower triangular, $U$ is upper triangular.\\
$A$ can reduced to echelon form without requiring row swaps:
\begin{enumerate}
    \item $A \xrightarrow{\text{row reduce}} U$.
    \item Find $L$: For each pivot col, divide the eliminated entries by the pivot for that col, and place it at the corresponding position.
    \begin{itemize}
        \item Alternative: $L = (E_p \dots E_1)^{-1}$, the elem matrices used to find $U$. Take the nonzero off-diagonal elements of the elementary inverse matrices and place them in the appropriate positions in $L$.
    \end{itemize}
\end{enumerate}
Then, $A\mathbf{x}=\mathbf{b} \to LU\mathbf{x}=\mathbf{b}$. Solve $\mathbf{y}$ where $L\mathbf{y}=\mathbf{b}$, then solve $U\mathbf{x}=\mathbf{y}$.\\
Zero pivot appears on the diagonal during row reduction $\Leftrightarrow$ row swap is needed.
\begin{itemize}
    \item Create permutation matrix $P$, which is $I$ with the required rows swapped.
    \item Then $PA=LU$.
\end{itemize}

\section*{Determinants}
\begin{enumerate}
    \item $|I|=1$.
    \item Swapping 2 rows/cols negates the det. Permutatn matrix with $r$ row exchanges: $|P| = (-1)^r$
    \item $\begin{vmatrix} ta & tb \\ c & d \end{vmatrix} = t \begin{vmatrix} a & b \\ c & d \end{vmatrix}$, $\quad$
    $\begin{vmatrix} a + a' & b + b' \\ c & d \end{vmatrix} = \begin{vmatrix} a & b \\ c & d \end{vmatrix} + \begin{vmatrix} a' + b' \\ c + d \end{vmatrix}$.
    \item 2 rows or cols are equal, then $|A| = 0$.
    \item If $\exists$ a zero row, then $|A| = 0$.
    \item Adding a scaled multiple of one row to another does not change determinant. So by Gaussian elimination, $|A| = \pm|U|$ (no scaling).\\
    $\begin{vmatrix} a & b \\ c - la & d - lb \end{vmatrix} = \begin{vmatrix} a & b \\ c & d \end{vmatrix}$.
    \item For a triangular matrix $A$, $|A| = a_{11}a_{22}\dots a_{nn}$.
    \item $|A| \neq 0 \Leftrightarrow A$ is invertible. Else $A$ is singular.
    \item $|AB| = |A||B|$ $\rightarrow$ $|A^{-1}| = \frac{1}{|A|}$. But $|A + B| \neq |A| + |B|$ generally.
    \item $|A^T| = |A|$.
\end{enumerate}
To find determinant: $A\xrightarrow{\text{reduce}}U$, by only swapping rows or add mult of one row to another. Then the determinant is the product of diagonals.\\
For $2\times 2$ or $3 \times 3$ matrix $A$, abs($|A|$) is the area of the parallelogram/parallelepiped formed by the cols, with one of the vertices at origin.\\
When applying linear transform $T$ determined by $A$ to a shape $S$, then $T(S) = \text{abs}(|A|) \times \text{original area of }S$.

\section*{Vector Space}
$V$ is nonempty set of vectors, closed under vector addition and scalar multiplication.
For $\mathbf{u},\mathbf{v},\mathbf{w},c,d$
\begin{enumerate}
    \item $\mathbf{0} \in V$.
    \item Addition: $\mathbf{u}+\mathbf{v} \in V$, $\mathbf{u}+\mathbf{v}=\mathbf{v}+\mathbf{u}$, $(\mathbf{u}+\mathbf{v})+\mathbf{w} = \mathbf{u}+(\mathbf{v}+\mathbf{w})$, $\exists -\mathbf{u}$ s.t. $\mathbf{u} + (-\mathbf{u}) = \mathbf{0}$.
    \item Multiplication: $c\mathbf{u}\in V$, $c(\mathbf{u} + \mathbf{v}) = c\mathbf{u} + c\mathbf{v}$, $(c + d)\mathbf{u} = c\mathbf{u} + d\mathbf{u}$, $c(d\mathbf{u}) = (cd)\mathbf{u}$, $1\mathbf{u} = \mathbf{u}$.
\end{enumerate}
Subspace $H$ is subset of $V$, that itself is a vector space.\\
Span of any set of vectors in $V$ is a subspace of $V$.\\
\textbf{Spanning Set Theorem}\\
If a vector in a spanning set is a linear combination of the others, you can safely remove it, and the remaining vectors will still span the space.
If $H \neq \{\mathbf{0}\}$, Some subset of the spanning set forms a basis for $H$.\\
A basis $\mathcal{B}$ for a subspace $H$ must be LI and a spanning set for $H$ ($H = Span\{\mathbf{b}_1,\ldots,\mathbf{b}_p\}$).\\
Cols of $I_n$ form a standard basis for $\mathbb{R}^n$. $\{1,t,\ldots,t^n\}$ forms basis for $\mathbb{P}_n$.\\
For any $m \times n$ matrix $A$,
\begin{itemize}
    \item Null space $\mathcal{N}(A)$: all solns $\mathbf{x}$ to $A\mathbf{x}=\mathbf{0}$. $\mathcal{N}(A)\subset\mathbb{R}^n$.
    \begin{itemize}
        \item To find basis: Row reducing $[A \mid \mathbf{0}]$, and express leading variables in terms of free variables.
        \item $\dim \mathcal{C}(A)$ is no. of pivot cols in $A$.
    \end{itemize}
    \item Col space $\mathcal{C}(A)$: Set of all linear combinations of cols of $A$. $\mathcal{C}(A)\subset\mathbb{R}^m$.
    \begin{itemize}
        \item To find basis: Row reduce $A$, identify cols with pivots. Extract the corresponding cols from \textbf{original} $A$. EROs change col space.
        \item $\dim \mathcal{N}(A)$ = rank($A$): no. of free variables in $A\mathbf{x}=\mathbf{0}$. Represent dim of set of $\mathbf{b}$ that make $A\mathbf{x}=\mathbf{b}$ consistent.
    \end{itemize}
    \item Row space $\mathcal{C}(A^T)$: Set of all linear combinations of rows of $A$. $\mathcal{C}(A^T)\subset \mathbb{R}^n$.
    \begin{itemize}
        \item To find basis: $A\xrightarrow{\text{REF}}B$. Nonzero rows of $B$ form basis. EROs do not change row space.
    \end{itemize}
\end{itemize}
R-N theorem: For matrix $A$ with $n$ cols: rank($A$) + $\dim \mathcal{N}(A) = n$.\\
\textbf{Coordinate systems}: If $\mathcal{B} = \{\mathbf{b}_1, \dots, \mathbf{b}_n\}$ is a basis for $V$, any vector $\mathbf{x}$ can be written uniquely as $\mathbf{x} = c_1\mathbf{b}_1 + \dots + c_n\mathbf{b}_n$. The column vector of these weights is the coordinate vector $[\mathbf{x}]_\mathcal{B}$.\\
The standard vector $\mathbf{x} = P_\mathcal{B}[\mathbf{x}]_\mathcal{B}$, where $P_\mathcal{B}$ contains the basis vectors as it cols.\\
To translate from $\mathcal{B}$ to $\mathcal{C}$: $[\mathbf{x}]_\mathcal{C} = P_{\mathcal{C} \leftarrow \mathcal{B}} [\mathbf{x}]_\mathcal{B}$.
Cols of $P_{\mathcal{C} \leftarrow \mathcal{B}}$ are simply the $\mathcal{B}$ basis vectors expressed in $\mathcal{C}$-coordinates: $[ [\mathbf{b}_1]_\mathcal{C} \dots [\mathbf{b}_n]_\mathcal{C} ]$. Set up $[\mathcal{C} \mid \mathcal{B}] \xrightarrow{\text{reduce}} [I \mid P_{\mathcal{C} \leftarrow \mathcal{B}}]$.


\section*{Orthogonality}
\textbf{Orthogonal} $\mathbf{u}$, $\mathbf{v}$: $\mathbf{u} \cdot \mathbf{v} = 0$; $||\mathbf{u}+\mathbf{v}||^2 = ||\mathbf{u}||^2 + ||\mathbf{v}||^2$.\\
\textbf{Orthonormal set}: Mutually orthogonal vectors, all norm = 1.\\

\textbf{Ortho comp} $W^{\perp} = \{\mathbf{z}:\mathbf{z}\cdot\mathbf{v}=0, \forall\mathbf{v}\in W\}$.\\
\begin{itemize}
    \item If $\mathbf{x}\in W$, $\mathbf{x} \perp$ to basis vectors of $W$.
\end{itemize}
\textbf{4 fundamental subspaces of a matrix}\\
$(Row A)^{\perp} = Nul A$ and $(Col A)^{\perp} = Nul A^T$.\\[-6pt]
\rule{\linewidth}{0.4pt}\\[-2pt]
\textbf{Ortho proj} of $\mathbf{u}$ onto reference vector $\mathbf{a}$: $proj_{\mathbf{a}}\mathbf{u} = \left(\frac{\mathbf{u} \cdot \mathbf{a}}{||\mathbf{a}||^2}\right)\mathbf{a}$.
Error vector $\mathbf{e} = \mathbf{u} - proj_{\mathbf{a}}\mathbf{u}$.\\
\textbf{Projection matrix} $P = \frac{\mathbf{v}\mathbf{v}^T}{\mathbf{v}^T \mathbf{v}}$ projects any vector in $\mathbb{R}^n$ onto $\mathbf{v}$.\\
\textbf{Best Approx Thm}: $||\mathbf{y} - \hat{\mathbf{y}}|| < ||\mathbf{y} - \mathbf{v}||$, $\forall \mathbf{v}\in W$.\\

% Ortho proj of vector $\mathbf{y}$ onto \textbf{subspace} $W$ with ortho basis $\{\mathbf{u}_1, \dots, \mathbf{u}_p\}$:
% $\hat{\mathbf{y}} = (\frac{\mathbf{y} \cdot \mathbf{u}_1}{\mathbf{u}_1 \cdot \mathbf{u}_1})\mathbf{u}_1 + \dots + (\frac{\mathbf{y} \cdot \mathbf{u}_p}{\mathbf{u}_p \cdot \mathbf{u}_p})\mathbf{u}_p$.
% If the basis is orthonormal, the denominators all become 1.\\

\noindent
\fbox{%
    \parbox{\dimexpr\linewidth-2\fboxsep-2\fboxrule\relax}{%
        Ortho proj of vector $\mathbf{y}$ onto \textbf{subspace} $W$ with ortho basis $\{\mathbf{u}_1, \dots, \mathbf{u}_p\}$:
        $\hat{\mathbf{y}} = \left(\frac{\mathbf{y} \cdot \mathbf{u}_1}{\mathbf{u}_1 \cdot \mathbf{u}_1}\right)\mathbf{u}_1 + \dots + \left(\frac{\mathbf{y} \cdot \mathbf{u}_p}{\mathbf{u}_p \cdot \mathbf{u}_p}\right)\mathbf{u}_p$.\\
        If the basis is orthonormal, the denominators all become 1.\\
        A non-zero set of ortho vectors is linearly dep.
    }%
}

\noindent
\fbox{%
    \parbox{\dimexpr\linewidth-2\fboxsep-2\fboxrule\relax}{%
        \textbf{Matrix $U$ w/ orthonormal cols}
        \begin{itemize}
            \item $U$ is either tall or square.
            \item For any $\mathbf{x}$, $\|U\mathbf{x}\|=\|\mathbf{x}\|$ and $(U\mathbf{x})\cdot(U\mathbf{y})=\mathbf{x}\cdot\mathbf{y}$
            \item $U^TU=I$ ($U^TU=D \iff U$ has \textbf{orthogonal} cols).
            \item \textbf{Projection matrix} $P = UU^T$: projects any vector $\mathbf{y}$ onto the col space of $U$ using $\hat{\mathbf{y}} = UU^T \mathbf{y}$.
            \begin{itemize}
                \item Rank($UU^T$) = Rank($U$) = Rank($U^T$) = Rank($U^TU$)
                \item $P^2 = P$ and $P^T = P$.
            \end{itemize}
            \item If $U$ is square:
            \begin{itemize}
                \item $U$ is orthogonal, $\det(U) = \pm 1$.
                \item $U^T U = UU^T = I$ and $U^{-1} = U^T$
            \end{itemize}
        \end{itemize}
    }%
}
\noindent
\fbox{%
    \parbox{\dimexpr\linewidth-2\fboxsep-2\fboxrule\relax}{%
        \textbf{QR-decomp}: $A=QR \implies A\mathbf{x}=\mathbf{y}\to R\mathbf{x}=Q^T\mathbf{y}$
        \begin{itemize}
            \item $A$ has linearly indep cols, $Q$ has orthonormal columns, $R$ is upper triangular, invertible.
        \end{itemize}
        \textbf{Gram-Schmidt process}:
        \begin{itemize}
            \item $\mathbf{v}_{k+1} = \mathbf{x}_{k+1} - \sum_{j=1}^{k} \text{proj}_{\mathbf{v}_j}\mathbf{x}_{k+1}$.
            \item $\mathbf{v}_1 = \mathbf{x}_1$, $\mathbf{v}_2 = \mathbf{x}_2 - \left(\frac{\mathbf{x}_2 \cdot \mathbf{v}_1}{\mathbf{v}_1 \cdot \mathbf{v}_1}\right)\mathbf{v}_1$, $\mathbf{v}_3 = \mathbf{x}_3 - \left(\frac{\mathbf{x}_3 \cdot \mathbf{v}_1}{\mathbf{v}_1 \cdot \mathbf{v}_1}\right)\mathbf{v}_1 - \left(\frac{\mathbf{x}_3 \cdot \mathbf{v}_2}{\mathbf{v}_2 \cdot \mathbf{v}_2}\right)\mathbf{v}_2$ $\ldots$
            \item Orthonormalise to form cols of $Q$: $\mathbf{u}_k = \frac{\mathbf{v}_k}{||\mathbf{v}_k||}$.
            \item $R = Q^T A$.
        \end{itemize}
        Economy QR: $m \times n$ $A$, $m \times n$ $Q$, $n \times n$ $R$.\\
        Full QR: $m \times m$ $Q$, $m \times n$ $R$.
        \begin{itemize}
            \item Append $I$ to $A$, form $Q_{\text{full}}$, then pad $R$ with rows of 0s.
        \end{itemize}
    }%
}

\section*{Least squares}
$A\mathbf{x} = \mathbf{b}$:
\begin{itemize}
    \item consistent: rank($A$)=rank($A|\mathbf{b}$)
    \item inconsistent: rank($A$) $<$ rank($A|\mathbf{b}$).
    $\mathbf{b}$ introduces new dimension, cannot be built from cols of $A$.
\end{itemize}
Usually $m > n$, overdetermined.\\[-6pt]
\rule{\linewidth}{0.4pt}\\[-3pt]
\textbf{Inconsistent}: error vector $\mathbf{e} = \mathbf{b} - A\mathbf{x}$.
\noindent
\fbox{%
    \parbox{\dimexpr\linewidth-2\fboxsep-2\fboxrule\relax}{%
        \textbf{Least squares goal}: minimise $||\mathbf{e}||^2$, i.e. $||\mathbf{b} - A\hat{\mathbf{x}}|| \le ||\mathbf{b} - A\mathbf{x}||$.\\
        \textbf{Method 1, Normal equation}: $A^T A\hat{\mathbf{x}} = A^T \mathbf{b}$. ($\because A^T\mathbf{e} = \mathbf{0}$)\\
        If cols of $A$ \textbf{linearly indep} $\to$ $A^TA$ invertible:
        \begin{itemize}
            \item \textbf{unique least-squares soln}: $\hat{\mathbf{x}} = (A^T A)^{-1}A^T \mathbf{b}$
            \item \textbf{Projectn matrix} $P = A(A^T A)^{-1}A^T$ takes vector $\mathbf{b}$ and outputs the closest point in the column space ($\hat{\mathbf{b}} = P\mathbf{b}$).
            \begin{itemize}
                \item Similar to $P=UU^T$ but cols of $A$ not orthonormal, only L.I.
            \end{itemize}
            % \item Projectn matrices $P^T = P$ (symmetric), $P^2 = P$ (idempotent).
        \end{itemize}
        \textbf{Method 2 QR decomp}: $R\hat{\mathbf{x}} = Q^T \mathbf{b}$.
    }%
}
\textbf{Applicatn}: $\mathbf{y} = X\boldsymbol{\beta} + \boldsymbol{\epsilon}$, $\mathbf{y}$ is the target data, $X$ is the design matrix, $\boldsymbol{\beta}$ contains the unknown weights, $\boldsymbol{\epsilon}$ is the residual error.\\
E.g. fit a line $y = \beta_0 + \beta_1 x$: 
\begin{itemize}
    \item $\mathbf{y}$ contains all y-coordinates.
    \item Design matrix $X$: Column 1 is all 1s (for $\beta_0$ intercept). Column 2 is all x-coordinates.
    \item Solve the Normal Equation $X^T X\boldsymbol{\beta} = X^T \mathbf{y}$ to find the optimal weights $\boldsymbol{\beta}$.
\end{itemize}
For polynomials, just add columns for $x^2$ etc. Linear wrt to weights not data.

\section*{Complex numbers}
$z = a + ib = \sqrt{a^2 + b^2} e^{i \theta} = |z|cis(\theta),\,\theta\in(-\pi,\pi]$\\
$i^1=i$, $i^2=-1$, $i^3=-i$, $i^4=1$, $\frac{1}{i} = -i$.\\
$\sqrt{-A}\sqrt{-B} = (i\sqrt{A})(i\sqrt{B})$.\\
\textbf{Inner product} $\langle \mathbf{a}, \mathbf{b} \rangle = \mathbf{b}^H \mathbf{a}$. $H$ means conjugate transpose.
\begin{itemize}
    \item $\mathbf{b}^H \mathbf{a} = \overline{\mathbf{a}^H \mathbf{b}} = (\mathbf{a}^H \mathbf{b})^H$.
\end{itemize}
\noindent
\fbox{%
    \parbox{\dimexpr\linewidth-2\fboxsep-2\fboxrule\relax}{%
        \textbf{DeMoivre Thm}: $z^n = r^n e^{in\theta} = r^n(\cos(n\theta) + i\sin(n\theta))$.\\
        \textbf{Applied to roots}: $z^{1/n} = r^{1/n}e^{i(\frac{\theta+2\pi k}{n})}$ for $k=0, 1, \dots, n-1$.
        \begin{itemize}
            \item $n$-th root will yield $n$ distinct roots, geometrically spaced out equally on a circle of radius $r^{1/n}$.
        \end{itemize}
    }%
}

\section*{DFT}
$\mathbf{e}_k[n] = e^{i\frac{2\pi}{N}kn}$ is a complex basis wave representing a specific frequency $k$ evaluated at time step $n$. +ve exponent means rotating anti-clockwise.
If the inner product of a signal vector and a basis wave is $0$, they are orthogonal, \& that specific frequency does not exist in the signal.\\
\textbf{DFT (Analysis)}: $\mathbf{X} = W\mathbf{x}$  or  $\mathbf{X}[k] = \sum_{n=0}^{N-1} \mathbf{x}[n]e^{-i\frac{2\pi}{N}kn}$.
\begin{itemize}
    \item Time-domain $\mathbf{x}$ (length $N$) decomposed to frequency-domain $\mathbf{X}$
    \item $W$ contains negative-exponents, is square, symmetric ($W^T = W$), ortho cols ($\frac{1}{N}W^TW = I$).
    \item Each row is one frequency ($k$) and $n$ increases with the columns.
\end{itemize}

% $$W = \begin{bmatrix}
% e^{-i\frac{2\pi}{N}(0)(0)} & e^{-i\frac{2\pi}{N}(0)(1)} & \dots & e^{-i\frac{2\pi}{N}(0)(N-1)} \\
% e^{-i\frac{2\pi}{N}(1)(0)} & e^{-i\frac{2\pi}{N}(1)(1)} & \dots & e^{-i\frac{2\pi}{N}(1)(N-1)} \\
% \vdots & \vdots & \ddots & \vdots \\
% e^{-i\frac{2\pi}{N}(N-1)(0)} & e^{-i\frac{2\pi}{N}(N-1)(1)} & \dots & e^{-i\frac{2\pi}{N}(N-1)(N-1)}
% \end{bmatrix}$$

% $$W =
% \begin{pmatrix}
% 1 & 1 & 1 & 1 \\
% 1 & e^{-i\pi/2} & e^{-i\pi} & e^{-i3\pi/2} \\
% 1 & e^{-i\pi} & e^{-i2\pi} & e^{-i3\pi} \\
% 1 & e^{-i3\pi/2} & e^{-i3\pi} & e^{-i9\pi/2}
% \end{pmatrix}
% $$

\vspace{0.5em}
\noindent
\resizebox{\columnwidth}{!}{%
$
\begin{bmatrix}
    X_0 \\
    \vdots \\
    \boxed{X_k} \\
    \vdots \\
    X_{N-1}
\end{bmatrix}
=
\left[
\begin{array}{cccc}
    e^{-i \frac{2\pi}{N} (0)(0)} & e^{-i \frac{2\pi}{N} (0)(1)} & \cdots & e^{-i \frac{2\pi}{N} (0)(N-1)} \\
    \vdots & \vdots & \ddots & \vdots \\
    \hline
    \multicolumn{1}{|c}{e^{-i \frac{2\pi}{N} (k)(0)}} & e^{-i \frac{2\pi}{N} (k)(1)} & \cdots & \multicolumn{1}{c|}{e^{-i \frac{2\pi}{N} (k)(N-1)}} \\
    \hline
    \vdots & \vdots & \ddots & \vdots \\
    e^{-i \frac{2\pi}{N} (N-1)(0)} & e^{-i \frac{2\pi}{N} (N-1)(1)} & \cdots & e^{-i \frac{2\pi}{N} (N-1)(N-1)}
\end{array}
\right]
\begin{bmatrix}
    \boxed{
    \begin{matrix}
        x_0 \\
        x_1 \\
        \vdots \\
        x_{N-1}
    \end{matrix}
    }
\end{bmatrix}
$
}

\textbf{4x4 fleshed out} $W_4 = \begin{bmatrix}
        1 & 1 & 1 & 1 \\
        1 & -i & -1 & i \\
        1 & -1 & 1 & -1 \\
        1 & i & -1 & -i
        \end{bmatrix}$
% \noindent
% \fbox{%
%     \parbox{\dimexpr\linewidth-2\fboxsep-2\fboxrule\relax}{%
%         % \textbf{4x4 DFT Matrix} ($W_4 = e^{-i 2\pi/4} = -i$):\\[-2pt]
%         % \rule{\linewidth}{0.4pt}\\[-2pt]
        
%     }%
% }

Each $\mathbf{X}[k]$ corresponds to freq $k$, and shows the magnitude and phase shift of that freq component in the signal $\mathbf{x}$.\\
For real-valued $\mathbf{x}$,
\begin{itemize}
    \item $\mathbf{X}[0]$ real.
    \item $\mathbf{X}[N/2]$ is Nyquist frequency, always real (for even $N$).
    \item $\mathbf{X}[N-k] = \overline{\mathbf{X}[k]}$: Rows $N/2 + 1$ to $N-1$ are just complex conjugates of the first half, redundant.
\end{itemize}

\textbf{Backward DFT (Synthesis)}: $\mathbf{x} = \frac{1}{N}W^H \mathbf{X}$ or $\mathbf{x}[n] = \frac{1}{N} \sum_{k=0}^{N-1} \mathbf{X}[k]e^{i\frac{2\pi}{N}kn}$.
\begin{itemize}
    \item $W^H$ contains positive exponents.
    \item Each row is one time ($n$) and $k$ increases with the columns.
\end{itemize}
% $$W^H = \begin{bmatrix}
% e^{i\frac{2\pi}{N}(0)(0)} & e^{i\frac{2\pi}{N}(1)(0)} & \dots & e^{i\frac{2\pi}{N}(N-1)(0)} \\
% e^{i\frac{2\pi}{N}(0)(1)} & e^{i\frac{2\pi}{N}(1)(1)} & \dots & e^{i\frac{2\pi}{N}(N-1)(1)} \\
% \vdots & \vdots & \ddots & \vdots \\
% e^{i\frac{2\pi}{N}(0)(N-1)} & e^{i\frac{2\pi}{N}(1)(N-1)} & \dots & e^{i\frac{2\pi}{N}(N-1)(N-1)}
% \end{bmatrix}$$


\noindent
\resizebox{\columnwidth}{!}{%
$
\begin{bmatrix}
    x_0 \\
    x_1 \\
    \vdots \\
    x_{N-1}
\end{bmatrix}
=
\frac{X_0}{N}
\begin{bmatrix}
    e^{i \frac{2\pi}{N} (0)(0)} \\
    e^{i \frac{2\pi}{N} (0)(1)} \\
    \vdots \\
    e^{i \frac{2\pi}{N} (0)(N-1)}
\end{bmatrix}
+ \cdots +
\frac{\boxed{X_k}}{N}
\begin{bmatrix}
\boxed{
\begin{matrix}
    e^{i \frac{2\pi}{N} (k)(0)} \\
    e^{i \frac{2\pi}{N} (k)(1)} \\
    \vdots \\
    e^{i \frac{2\pi}{N} (k)(N-1)}
\end{matrix}
}
\end{bmatrix}
+ \cdots +
\frac{X_{N-1}}{N}
\begin{bmatrix}
    e^{i \frac{2\pi}{N} (N-1)(0)} \\
    e^{i \frac{2\pi}{N} (N-1)(1)} \\
    \vdots \\
    e^{i \frac{2\pi}{N} (N-1)(N-1)}
\end{bmatrix}
$
}

\section*{Eigenvectors}
$n \times n$ $A\mathbf{x} = \lambda\mathbf{x}$, $\mathbf{x}\neq \mathbf{0}$.\\
$\lambda = 0 \iff A\mathbf{x} = \mathbf{0}$, $\mathbf{x}\neq \mathbf{0} \iff A$ not invertible.\\
If $\lambda$ is eigenvalue of $A$, and all $\lambda \neq 0$, then $\frac{1}{\lambda}$ is eigenvalue of $A^{-1}$.\\
$n \times n$ $A$ has $n$ eigenvalues, counting algebraic multiplicities and allowed to be complex.\\
$\lambda_1 + \ldots + \lambda_n = $Tr$(A)$ and $\lambda_1\lambda_2\ldots\lambda_n = \det(A)$\\
$2 \times 2$ $A = \begin{bmatrix} a & b \\ c & d \end{bmatrix}$, $\det(A-\lambda I) = \lambda^2 - \text{Tr}(A)\lambda + \det(A) = 0$
\noindent
\fbox{%
    \parbox{\dimexpr\linewidth-2\fboxsep-2\fboxrule\relax}{%
        \textbf{Finding} $\lambda$: 
        \begin{itemize}
            \item Diagonal/triangular matrix: eigenvalues = main diagonal entries.
            \item Characteristic eqn: $(A - \lambda I)\mathbf{x} = \mathbf{0} \rightarrow \det(A - \lambda I) = 0$.
            The multiplicity of a root $\lambda$ is its \textbf{algebraic multiplicity}.
        \end{itemize}
        \textbf{Finding eigenvectors}:
        \begin{itemize}
            \item Row reduce $(A - \lambda I)\mathbf{x} = \mathbf{0}$ to get general solution for $\mathbf{x}$.
            \item The basis vectors of this null space are the eigenvectors for that eigenspace. The number of basis vectors is the \textbf{geometric mult}.
            \item \textbf{Defective matrix}: If GM $<$ AM for any $\lambda$, the matrix is defective (cannot be diagonalized).
        \end{itemize}
    }%
}

% \noindent
% \fbox{%
%     \parbox{\dimexpr\linewidth-2\fboxsep-2\fboxrule\relax}{%
%         \textbf{Finding} $\lambda$: 
%         \begin{itemize}
%             \item Diagonal/triangular matrix: eigenvalues = main diag entries.
%             \item \textbf{Characteristic eqn}: $\det(A - \lambda I) = 0$. 
            % \item $2 \times 2$ shortcut: $\lambda^2 - \text{tr}(A)\lambda + \det(A) = 0$.
            % \item Properties: $\sum \lambda_i = \text{tr}(A)$ and $\prod \lambda_i = \det(A)$.
%         \end{itemize}
%     }%
% }

\textbf{Similar matrices} $B = P^{-1}AP$ represent the same linear transformation in different coordinate systems. They have the same determinant, invertibility, rank, nullity, trace, characteristic polynomial, and eigenvalues.\\
\noindent
\fbox{%
    \parbox{\dimexpr\linewidth-2\fboxsep-2\fboxrule\relax}{%
        \textbf{Diagonalisation Thm}: $A = PDP^{-1}$. 
        \begin{itemize}
            \item If an $n \times n$ matrix has $n$ unique eigenvalues, it is diagonalisable.
            \begin{itemize}
                \item Eigenvectors from diff eigenvalues are linearly independent.
            \end{itemize}
            \item Otherwise, it must have $n$ linearly indep eigenvectors (enough to form a full basis for $\mathbb{R}^n$; i.e. an \textbf{eigenbasis}).
            \begin{itemize}
                \item For repeated eigenvalues, their \textbf{AM} must equal \textbf{GM}.
            \end{itemize}
            \item Cols of $P$ are the eigenvectors. $D$ is diagonal matrix w/ the corresponding eigenvalues. $\lambda_k$ in $D$ must match $\mathbf{v}_k$ in $P$.
        \end{itemize}
    }%
}
\textbf{Invertibility vs diagonalisability}
I \& D: $\begin{bmatrix} 2 & 0 \\ 0 & 3 \end{bmatrix}$;
not I \& D: $\begin{bmatrix} 1 & 0 \\ 0 & 0 \end{bmatrix}$;
I \& not D: $\begin{bmatrix} 1 & 1 \\ 0 & 1 \end{bmatrix}$;
not I \& not D: $\begin{bmatrix} 0 & 1 \\ 0 & 0 \end{bmatrix}$

\textbf{Higher powers of $A$}: $A^k = PD^kP^{-1}$\\
\textbf{Decoupling}: $A\mathbf{x} = PDP^{-1}\mathbf{x}$ first translates $\mathbf{x}$ to eigenbasis ($P^{-1})$, transforms it by the eigenvalues ($D$), and translates back to standard basis ($P$).
Recall that $[x]_B = P_B^{-1}x$.\\
% \textbf{Dynamic system $x_k = A^k x_0$}\\
% With an eigenbasis, any initial vector can be written as
% $\mathbf{x}_0 = c_1\mathbf{v}_1 + \ldots + c_n\mathbf{v}_n$\\
% $x_k = c_1\lambda^k\mathbf{v}_1 + \ldots + c_n\lambda_n^k \mathbf{v}_n$
% \begin{itemize}
%     \item Orthogonal eigenvectors: $c_i = \frac{\mathbf{x}_0\cdot \mathbf{v}_i}{\mathbf{v}_i\cdot\mathbf{v}_i}$
%     \item Non-ortho eigenvectors: solve $Pc = \mathbf{x}_0$
% \end{itemize}

\noindent
\fbox{%
    \parbox{\dimexpr\linewidth-2\fboxsep-2\fboxrule\relax}{%
        \textbf{Dynamic system $\mathbf{x}_k = A^k \mathbf{x}_0 = PD^kP^{-1}\mathbf{x}_0$}\\
        With an eigenbasis, any initial vector can be written as:\\
        $\mathbf{x}_0 = c_1\mathbf{v}_1 + \ldots + c_n\mathbf{v}_n$\\
        $\mathbf{x}_k = c_1(\lambda_1)^k\mathbf{v}_1 + \ldots + c_n(\lambda_n)^k \mathbf{v}_n$
        \begin{itemize}
            \item Ortho eigenvectors: $c_i = \frac{\mathbf{x}_0\cdot \mathbf{v}_i}{\mathbf{v}_i\cdot\mathbf{v}_i}$.
            \item Non-ortho eigenvectors: solve $P\mathbf{c} = \mathbf{x}_0$ (cols of $P$ are $\mathbf{v}_i$).
            \item \textbf{Long-term behavior} ($k \to \infty$): 
            \begin{itemize}
                \item Terms with $|\lambda_i| < 1$ shrink to $\mathbf{0}$.
                \item System is pulled toward the \textbf{dominant eigenvector} (the one with the largest $|\lambda|$).
            \end{itemize}
        \end{itemize}
    }%
}

\newpage

\section*{Extras}
\textbf{Extra visualisation of DFT synthesis}
\noindent
\resizebox{\columnwidth}{!}{%
$
\begin{bmatrix}
    x_0 \\
    \vdots \\
    \boxed{x_n} \\
    \vdots \\
    x_{N-1}
\end{bmatrix}
=
\frac{1}{N}
\left[
\begin{array}{cccc}
    e^{i \frac{2\pi}{N} (0)(0)} & e^{i \frac{2\pi}{N} (0)(1)} & \cdots & e^{i \frac{2\pi}{N} (0)(N-1)} \\
    \vdots & \vdots & \ddots & \vdots \\
    \hline
    \multicolumn{1}{|c}{e^{i \frac{2\pi}{N} (n)(0)}} & e^{i \frac{2\pi}{N} (n)(1)} & \cdots & \multicolumn{1}{c|}{e^{i \frac{2\pi}{N} (n)(N-1)}} \\
    \hline
    \vdots & \vdots & \ddots & \vdots \\
    e^{i \frac{2\pi}{N} (N-1)(0)} & e^{i \frac{2\pi}{N} (N-1)(1)} & \cdots & e^{i \frac{2\pi}{N} (N-1)(N-1)}
\end{array}
\right]
\begin{bmatrix}
    \boxed{
    \begin{matrix}
        X_0 \\
        X_1 \\
        \vdots \\
        X_{N-1}
    \end{matrix}
    }
\end{bmatrix}
$
}

\vspace{0.5em}
\textbf{Fleshed out DFT matrices}\\
\vspace{2pt}
\noindent
\fbox{%
    \parbox{\dimexpr\linewidth-2\fboxsep-2\fboxrule\relax}{%
        \textbf{8x8 DFT Matrix} ($W = e^{-i 2\pi/8}$):\\[-2pt]
        \rule{\linewidth}{0.4pt}\\[-2pt]
        {\setlength{\arraycolsep}{3pt}\scriptsize
        $W_8 = \begin{bmatrix}
        1 & 1 & 1 & 1 & 1 & 1 & 1 & 1 \\
        1 & W & W^2 & W^3 & W^4 & W^5 & W^6 & W^7 \\
        1 & W^2 & W^4 & W^6 & 1 & W^2 & W^4 & W^6 \\
        1 & W^3 & W^6 & W & W^4 & W^7 & W^2 & W^5 \\
        1 & W^4 & 1 & W^4 & 1 & W^4 & 1 & W^4 \\
        1 & W^5 & W^2 & W^7 & W^4 & W & W^6 & W^3 \\
        1 & W^6 & W^4 & W^2 & 1 & W^6 & W^4 & W^2 \\
        1 & W^7 & W^6 & W^5 & W^4 & W^3 & W^2 & W
        \end{bmatrix}$}
    }%
}

\noindent
\fbox{%
    \parbox{\dimexpr\linewidth-2\fboxsep-2\fboxrule\relax}{%
        \textbf{8x8 DFT Matrix} ($W = e^{-i 2\pi/8}$):\\[-2pt]
        Let $a = \frac{1}{\sqrt{2}}$. Note $W^1 = a - ia$, $W^2 = -i$.\\[-2pt]
        \rule{\linewidth}{0.4pt}\\[-2pt]
        {\setlength{\arraycolsep}{2pt}\tiny
        $W_8 = \begin{bmatrix}
        1 & 1 & 1 & 1 & 1 & 1 & 1 & 1 \\
        1 & a-ia & -i & -a-ia & -1 & -a+ia & i & a+ia \\
        1 & -i & -1 & i & 1 & -i & -1 & i \\
        1 & -a-ia & i & a-ia & -1 & a+ia & -i & -a+ia \\
        1 & -1 & 1 & -1 & 1 & -1 & 1 & -1 \\
        1 & -a+ia & -i & a+ia & -1 & a-ia & i & -a-ia \\
        1 & i & -1 & -i & 1 & i & -1 & -i \\
        1 & a+ia & i & -a+ia & -1 & -a-ia & -i & a-ia
        \end{bmatrix}$}
    }%
}

\vspace{0.5em}
\textbf{Extra stuff about eigenvalues}
\noindent
\fbox{%
    \parbox{\dimexpr\linewidth-2\fboxsep-2\fboxrule\relax}{%
        \textbf{Spectral Theorem} (For Symmetric Matrices $A = A^T$):
        \begin{itemize}
            \item All eigenvalues ($\lambda$) are real numbers.
            \item Eigenvectors from different $\lambda$ are automatically \textbf{mutually orthogonal} ($\mathbf{v}_i \cdot \mathbf{v}_j = 0$).
            \item $A$ is always orthogonally diagonalisable: $A = QDQ^T$.
            \begin{itemize}
                \item $Q$ is an orthogonal matrix (cols are orthonormal eigenvectors).
                \item Since $Q$ is orthogonal, $Q^{-1} = Q^T$.
            \end{itemize}
            \item \textbf{Spectral Decomposition}: $A = \lambda_1\mathbf{u}_1\mathbf{u}_1^T + \dots + \lambda_n\mathbf{u}_n\mathbf{u}_n^T$ (where $\mathbf{u}_i$ are the orthonormal columns of $Q$).
        \end{itemize}
    }%
}

\newpage


\end{multicols*}
\end{document}